
\documentclass[11pt,a4paper]{scrartcl}

% ==========================================================
% Pakete
% ==========================================================
\usepackage[T1]{fontenc}
\usepackage[ngerman]{babel}
\usepackage{lmodern}
\usepackage{mathtools}
\usepackage{amssymb}
\usepackage{cancel} %für mathematische Ausdrücke durchzustreichen oder zu kürzen

% ==========================================================
% Eigene mathematische Operatoren
% ==========================================================
\DeclareMathOperator{\Var}{Var}
\DeclareMathOperator{\Cov}{Cov}
\DeclareMathOperator{\diag}{diag}

\DeclarePairedDelimiter{\abs}{\lvert}{\rvert}
\DeclarePairedDelimiter{\norm}{\lVert}{\rVert}

\begin{document}
	
	\section{Mathematische Befehle und Symbole}
	\label{sec:mathematische-befehle}
	
	Dieses Kapitel enthält eine Sammlung der wichtigsten
	mathematischen Befehle in LaTeX.
	
	% ==========================================================
	\subsection{Brüche, Potenzen und Wurzeln}
	% ==========================================================
	
	\begin{align}
		\text{Bruch:}\quad %erzeugt einen horizontalen Abstand
		&\frac{a}{b} \\[4pt]
		\text{Großer Bruch:}\quad
		&\dfrac{a+b}{c+d} \\[4pt]
		\text{Kleiner Bruch:}\quad
		&\tfrac{a}{b} \\[4pt]
		\text{Quadratwurzel:}\quad
		&\sqrt{x} \\[4pt]
		\text{n-te Wurzel:}\quad
		&\sqrt[n]{x} \\[4pt]
		\text{Potenz:}\quad
		&x^{n+1} \\[4pt]
		\text{Index:}\quad
		&x_{i,j} \\[4pt]
		\text{Binomialkoeffizient:}\quad
		&\binom{n}{k}
	\end{align}
\\


	
	\[ % keine Nummerierung
	\begin{alignedat}{2}
		\text{Bruch:}          &\qquad& \frac{a}{b} \\
		\text{Großer Bruch:}   &      & \dfrac{a+b}{c+d} \\
		\text{Kleiner Bruch:}  &      & \tfrac{a}{b} \\
		\text{Quadratwurzel:} &      & \sqrt{x}
	\end{alignedat}
	\]
\\

	
	\begin{equation} % wird nummeriert
		\begin{alignedat}{1}
			\text{Bruch:}\qquad          & \frac{a}{b} \\
			\text{Großer Bruch:}\qquad   & \dfrac{a+b}{c+d} \\
			\text{Kleiner Bruch:}\qquad  & \tfrac{a}{b} \\
			\text{Quadratwurzel:}\qquad & \sqrt{x}
		\end{alignedat}
		\label{eq:formelsammlung}
	\end{equation}
\\
	
	Beispiel: Quadratische Lösungsformel
	
	\begin{equation}
		\boxed{
			x_{1,2} =
			\frac{-b \pm \sqrt{b^2-4ac}}{2a}
		}
	\end{equation}
	
	% ==========================================================
	\subsection{Summen, Produkte und Grenzwerte}
	% ==========================================================
	
	\begin{align}
		\sum_{i=1}^{n} i
		&= \frac{n(n+1)}{2} \\[6pt]
		\prod_{i=1}^{n} i
		&= n! \\[6pt]
		\lim_{x\to 0}\frac{\sin x}{x}
		&= 1 \\[6pt]
		\lim_{n\to\infty}\frac{1}{n}
		&= 0 \\[6pt]
		\limsup_{n\to\infty}x_n
		& \\[6pt]
		\liminf_{n\to\infty}x_n
		&
	\end{align}
	
	% ==========================================================
	\subsection{Integrale}
	% ==========================================================
	
	\begin{align}
		\int f(x)\,dx
		&= F(x)+C \\[6pt]
		\int_a^b f(x)\,dx
		&= F(b)-F(a) \\[6pt]
		\iint_D f(x,y)\,dA
		& \\[6pt]
		\iiint_V f(x,y,z)\,dV
		& \\[6pt]
		\oint_C \vec{F}\cdot d\vec{r}
		&
	\end{align}
	
	% ==========================================================
	\subsection{Ableitungen}
	% ==========================================================
	
	\begin{align}
		f'(x)
		&= \frac{df}{dx} \\[6pt]
		f''(x)
		&= \frac{d^2f}{dx^2} \\[6pt]
		\frac{\partial f}{\partial x}
		& \\[6pt]
		\frac{\partial^2 f}{\partial x^2}
		& \\[6pt]
		\dot{x}
		&= \frac{dx}{dt} \\[6pt]
		\ddot{x}
		&= \frac{d^2x}{dt^2}
	\end{align}
	
	% ==========================================================
	\subsection{Mathematische Funktionen}
	% ==========================================================
	
	\begin{align}
		\sin(x) &\qquad \cos(x) \qquad \tan(x)\\
		\arcsin(x) &\qquad \arccos(x)\\
		\ln(x) &\qquad \log(x)\\
		\log_2(x) &\qquad \exp(x)\\
		\min_{x\in A}f(x)
		&\qquad \max_{x\in A}f(x)
	\end{align}
	
	Beispiel: Exponentialfunktion
	
	\begin{equation}
		f(x)=A\exp(-\lambda x)
	\end{equation}
	
	% ==========================================================
	\subsection{Vektoren und Matrizen}
	% ==========================================================
	
	\begin{align}
		\vec{v} &=
		\begin{pmatrix}
			v_x\\
			v_y\\
			v_z
		\end{pmatrix}
		\\[8pt]
		A &=
		\begin{pmatrix}
			a & b\\
			c & d
		\end{pmatrix}
		\\[8pt]
		B &=
		\begin{bmatrix}
			1 & 2\\
			3 & 4
		\end{bmatrix}
	\end{align}
	
	Determinante:
	
	\begin{equation}
		\det(A)=
		\begin{vmatrix}
			a & b\\
			c & d
		\end{vmatrix}
		=ad-bc
	\end{equation}
	
	Vektornorm:
	
	\begin{equation}
		\norm{\vec{v}}
		=\sqrt{v_x^2+v_y^2+v_z^2}
	\end{equation}
	
	% ==========================================================
	\subsection{Mengen und logische Operatoren}
	% ==========================================================
	
	\begin{align*}
		a &\neq b\\
		a &\leq b\\
		a &\geq b\\
		a &\approx b\\
		a &\equiv b\\
		a &\propto b\\
		x &\in A\\
		x &\notin A\\
		A &\subseteq B\\
		A &\cup B\\
		A &\cap B
	\end{align*}
	
	Weitere Symbole:
	
	\[
	\forall,\quad
	\exists,\quad
	\emptyset,\quad
	\land,\quad
	\lor,\quad
	\Rightarrow,\quad
	\Leftrightarrow
	\]
	
	% ==========================================================
	\subsection{Griechische Buchstaben}
	% ==========================================================
	
	\[
	\begin{gathered}
		\alpha,\ \beta,\ \gamma,\ \delta,\ \epsilon,\
		\varepsilon,\ \zeta,\ \eta,\ \theta,\ \vartheta,\\
		\iota,\ \kappa,\ \lambda,\ \mu,\ \nu,\
		\xi,\ \pi,\ \rho,\ \sigma,\ \tau,\\
		\upsilon,\ \phi,\ \varphi,\ \chi,\ \psi,\ \omega
	\end{gathered}
	\]
	
	Großbuchstaben:
	
	\[
	\Gamma,\ \Delta,\ \Theta,\ \Lambda,\
	\Xi,\ \Pi,\ \Sigma,\ \Upsilon,\
	\Phi,\ \Psi,\ \Omega
	\]
	
	% ==========================================================
	\subsection{Klammern und Beträge}
	% ==========================================================
	
	Automatisch angepasste Klammergrößen:
	
	\begin{equation}
		f(x)=
		\left(
		\frac{x^2+1}{x-1}
		\right)^2
	\end{equation}
	
	Betrag:
	
	\begin{equation}
		\abs{x},\qquad
		\abs*{\frac{x+1}{x-1}}
	\end{equation}
	
	% ==========================================================
	\subsection{Weitere nützliche Befehle}
	% ==========================================================
	
	\begin{align}
		a &\overset{!}{=} b\\
		x &\xrightarrow{n\to\infty} 0\\
		f(x) &\coloneqq x^2+3x+2\\
		\overline{x} &= \frac{1}{n}
		\sum_{i=1}^{n}x_i\\
		\hat{\theta} &= \theta+\varepsilon\\
		\cancel{x} + y &= y
		\quad \text{für }x=0
	\end{align}
	
	Terme beschriften:
	
	\begin{equation}
		F =
		\underbrace{ma}_{\text{Trägheit}}
		+
		\underbrace{c\dot{x}}_{\text{Dämpfung}}
		+
		\underbrace{kx}_{\text{Federkraft}}
	\end{equation}
	
	% ==========================================================
	\subsection{Eigene mathematische Operatoren}
	% ==========================================================
	
	\begin{equation}
		\Var(X+Y)
		=
		\Var(X)+\Var(Y)+2\Cov(X,Y)
	\end{equation}
	
	\begin{equation}
		D=\diag(\lambda_1,\lambda_2,\lambda_3)
	\end{equation}
	
	% ==========================================================
	\subsection{Fallunterscheidungen}
	% ==========================================================
	
	\begin{equation}
		f(x)=
		\begin{dcases}
			x^2, & \text{wenn }x>0,\\
			0, & \text{wenn }x=0,\\
			-x, & \text{wenn }x<0.
		\end{dcases}
	\end{equation}
	
	% ==========================================================
	\subsection{Gleichungssysteme}
	% ==========================================================
	
	\begin{equation}
		\begin{aligned}
			2x+y &= 5\\
			x-y &= 1
		\end{aligned}
	\end{equation}
	
	% ==========================================================
	\subsection{Lange Gleichungen}
	% ==========================================================
	
	\begin{align}
		\MoveEqLeft
		Q_0=-m_0gl_0\sin(\theta_0)
		\notag\\
		&+m_1gL_1\cos(\theta_0)
		\notag\\
		&-m_2gL_2\cos(\theta_0)
		\notag\\
		&-D_0\dot{\theta}_0
		\notag\\
		&+D_1(\dot{\theta}_1-\dot{\theta}_0)
		\notag\\
		&+D_2(\dot{\theta}_2-\dot{\theta}_0)
		\notag\\
		&+M_{\mathrm{control}}
	\end{align}

	\subsection{Brüche kürzen}
% ==========================================================

	\begin{equation}
		\frac{2\cancel{x}y}{4\cancel{x}}
		= \frac{y}{2}
	\end{equation}

	\begin{equation}
		a+b-\cancelto{0}{b-b}=a
	\end{equation}
	
	\begin{equation}
		a+\cancelto{0}{(b-b)}=a
	\end{equation}

	\begin{equation}
		\lim_{x\to 0}
		\frac{x^2+3x}{x}
		=
		\lim_{x\to 0}
		\frac{\cancel{x}(x+3)}{\cancel{x}}
		=3
	\end{equation}

\end{document}
